% The abstract. `make paper-text` also writes it as plain text to arxiv/abstract.txt (180 to 250 words).
Can AI models that read financial news predict how stock prices will react, or tell when the market has over- or
under-reacted? We test this with AI decision models, non-generative language models that return probabilities for
typed questions, in a pre-registered, look-ahead-free design on S\&P 500 stocks. We fine-tune an open-weights
decision model on 2024--2025 news and returns, compare it with commercial decision models and standard baselines,
mask company identities, freeze the analysis plan with a cryptographic hash, and unblind a January--July 2026
test period once. The fine-tuned model's signal replicated on two validation periods (daily information
coefficient \res{lf-ic-valb}, $t = \res{lf-ic-valb-t1}$) but fell to \res{lf-ic-test} ($t = \res{lf-ic-test-t1}$)
on the blinded test; none of the five primary hypotheses is supported. Where the models predict anything, it is
the overnight gap: commercial models forecast the move to the open (information coefficient \res{d1-ic-on-test}
to \res{drex-ic-on-test}) but not the move after it. The models cannot tell over- from under-reaction, and no
pre-registered trading strategy survives transaction costs and a deflated Sharpe ratio. Commercial models
identify masked companies in \res{mem-min} to \res{mem-max} of cases against a \res{mem-baseline} baseline, a
warning about look-ahead bias in backtests of large language models (LLMs) on historical news. In an exploratory
analysis, conditions selected on two validation periods, with in-sample Sharpe ratios up to
\res{cb-valb-gross-max}, failed on the test. The market reads the news first.
